% =======================================================================
% Relativistic Extension of Chapter X, Sections 96--97
% Magnetic Field Effects on Rayleigh--Taylor Instability
%
% Conventions: (-,+,+,+) signature, c explicit, u^mu u_mu = -c^2
% Requires: rel_preamble.tex macros
% =======================================================================

\chapter{Relativistic Rayleigh--Taylor Instability with Magnetic Fields
  (Ch.~X, \S\S96--97)}\label{ch:rel-10-96-97}

In Chapter~X, \S\S96--97, Chandrasekhar analyses the stabilising
influence of an impressed magnetic field on the Rayleigh--Taylor (RT)
instability of two superposed fluids.  When the magnetic field is strong
enough that the Alfv\'en speed~$\vA$ becomes an appreciable fraction of
the speed of light, the classical analysis must be replaced by a fully
relativistic treatment.  This chapter develops that extension for both
vertical (\S\ref{sec:rel-96}) and horizontal
(\S\ref{sec:rel-97}) magnetic fields.  The speed of light~$c$ is kept
explicit throughout, so that every classical limit is recovered
transparently by $c\to\infty$.

%======================================================================
\section{Relativistic preliminaries for magnetised RT
  instability}\label{sec:rel-96-prelim}
%======================================================================

\subsection{Relativistic enthalpy and Alfv\'en speed}

In the relativistic regime the inertial role of the mass density~$\rho$
is assumed by the relativistic enthalpy density
\begin{equation}
  \enthalpy = \varepsilon + p,
  \label{eq:rel96-enthalpy}
\end{equation}
where $\varepsilon$ is the total energy density (including rest-mass
energy $\rho c^{2}$) and $p$ the isotropic pressure.  The effective
inertia per unit volume appearing in the equation of motion is
$\enthalpy/c^{2}$.

The magnetic field in the fluid rest frame is described by the 4-vector
$b^{\mu}$, with invariant magnitude $b^{2}=b^{\mu}b_{\mu}$.  The
relativistic Alfv\'en speed along the background field direction is
(cf.~RELATIVISTIC\_CONVENTIONS.md)
\begin{equation}
  \vA^{2}
    = \frac{b^{2}\,c^{2}}{\enthalpy + b^{2}}
    = \frac{b^{2}\,c^{2}}{4\pi\enthalpy/\mu + b^{2}}\;\Big|_{\text{Gaussian}},
  \label{eq:rel96-vA}
\end{equation}
where in the Gaussian form $b^{2}=\mu H^{2}/(4\pi)$.

\paragraph{Causality constraint.}
Because the denominator $\enthalpy + b^{2}>0$, one has
$\vA^{2}<c^{2}$ for any finite magnetic field and positive enthalpy.
Thus \emph{the Alfv\'en speed is automatically sub-luminal}, ensuring
that magnetic stabilisation of the RT instability never violates
causality.

\subsection{Classical limit}

In the non-relativistic limit $\enthalpy\to\rho c^{2}$ and
$b^{2}\ll\rho c^{2}$, so that
\begin{equation}
  \vA^{2}\;\longrightarrow\;
  \frac{\mu H^{2}}{4\pi\rho}
  = V_{a}^{2}\quad(c\to\infty),
  \label{eq:rel96-vA-NR}
\end{equation}
recovering the classical Alfv\'en speed of Chapter~IV, \S\,38.

\subsection{Relativistic ideal-MHD energy-momentum tensor}

The total energy-momentum tensor for a perfect, magnetised fluid is
\begin{equation}
  T^{\mu\nu}
    = \Bigl(\enthalpy + b^{2}\Bigr)\frac{u^{\mu}u^{\nu}}{c^{2}}
      + \Bigl(p + \tfrac{1}{2}b^{2}\Bigr)g^{\mu\nu}
      - b^{\mu}b^{\nu}.
  \label{eq:rel96-Tmunu}
\end{equation}
This furnishes the starting point for the linearised perturbation
analysis that follows.

%======================================================================
\section{The effect of a vertical magnetic field --- relativistic
  extension}\label{sec:rel-96}
%======================================================================

We now extend \S96 of Chandrasekhar to the relativistic regime.  A
uniform magnetic field $\mathbf{B}=B\hat{\mathbf{z}}$ is impressed
parallel to the gravitational acceleration
$\mathbf{g}=-g\hat{\mathbf{z}}$.  Two uniform fluids of relativistic
enthalpy densities $\enthalpy_{1}$ (below, $z<0$) and $\enthalpy_{2}$
(above, $z>0$) are separated by a horizontal interface at $z=0$.

\subsection{Linearised relativistic perturbation equations}

Linearising the conservation law $\nabla_{\mu}T^{\mu\nu}=0$ about a
static equilibrium with 4-velocity $u^{\mu}=(c,\mathbf{0})$ and a
uniform vertical field $b^{\mu}=(0,0,0,B_{0})$ in the fluid frame, we
obtain the relativistic counterparts of Chandrasekhar's
equations~\eqref{eq:10-171}--\eqref{eq:10-174}.  For perturbations with
the dependence $\exp\!\bigl[i(k_{x}x+k_{y}y)+nt\bigr]$:
\begin{gather}
  \frac{\enthalpy}{c^{2}}\,n\,u
    - \frac{b^{2}}{c^{2}}\,\frac{1}{n}\bigl(D h_{x}-ik_{x}h_{z}\bigr)
    = -ik_{x}\,\delta p,
  \label{eq:rel96-perturb-x}\\
  \frac{\enthalpy}{c^{2}}\,n\,v
    - \frac{b^{2}}{c^{2}}\,\frac{1}{n}\bigl(D h_{y}-ik_{y}h_{z}\bigr)
    = -ik_{y}\,\delta p,
  \label{eq:rel96-perturb-y}\\
  \frac{\enthalpy}{c^{2}}\,n\,w
    = -D\,\delta p + \frac{g}{n}\,D\!\Bigl(\frac{\enthalpy}{c^{2}}\Bigr)\,w,
  \label{eq:rel96-perturb-z}\\
  h_{z} = \frac{B_{0}}{n}\,Dw,
  \label{eq:rel96-hz}
\end{gather}
together with the solenoidal conditions $ik_{x}u+ik_{y}v+Dw=0$ and
$ik_{x}h_{x}+ik_{y}h_{y}+Dh_{z}=0$.  Here $D=d/dz$, and we note that
the classical density~$\rho$ is everywhere replaced by the relativistic
inertia $\enthalpy/c^{2}$.

Following the same algebraic steps as in
\S96 (eliminating $\delta p$ between the horizontal and vertical
equations), one arrives at the relativistic master equation:
\begin{equation}
  D\!\Bigl(\frac{\enthalpy}{c^{2}}\,Dw\Bigr)
    - \frac{\vA^{2}}{n^{2}}\,\frac{\enthalpy}{c^{2}\!-\!\vA^{2}}\,
      (D^{2}-k^{2})\,D^{2}w
    = k^{2}\frac{\enthalpy}{c^{2}}\,w
      - \frac{gk^{2}}{n^{2}}\,D\!\Bigl(\frac{\enthalpy}{c^{2}}\Bigr)\,w.
  \label{eq:rel96-master}
\end{equation}
In the limit $c\to\infty$ (so that $\vA^{2}/c^{2}\to0$ and
$\enthalpy/c^{2}\to\rho$), the factor
$\vA^{2}\enthalpy/(c^{2}-\vA^{2})\to\mu H^{2}/(4\pi)$ and
equation~\eqref{eq:rel96-master} reduces to Chandrasekhar's
equation~\eqref{eq:10-183}.

\subsection{Integral relation and reality of $n^{2}$}

Multiplying~\eqref{eq:rel96-master} by $w^{*}$ and integrating between
the bounding surfaces (where $w=0$), we obtain
\begin{equation}
  \int\frac{\enthalpy}{c^{2}}
    \bigl\{|Dw|^{2}+k^{2}|w|^{2}\bigr\}\,dz
  + \frac{\vA^{2}}{n^{2}}\int\frac{\enthalpy}{c^{2}\!-\!\vA^{2}}
    \bigl\{|D^{2}w|^{2}+k^{2}|Dw|^{2}\bigr\}\,dz
  = \frac{gk^{2}}{n^{2}}\int D\!\Bigl(\frac{\enthalpy}{c^{2}}\Bigr)
    |w|^{2}\,dz.
  \label{eq:rel96-integral}
\end{equation}
Since $\vA^{2}<c^{2}$ (causality), all coefficients in the integrands
are positive-definite.  The argument of
Chandrasekhar (cf.\ equation~\eqref{eq:10-185}) therefore carries over:
\emph{the characteristic values of $n^{2}$ are necessarily real in the
relativistic theory as well}.

\subsection{Two uniform fluids: the unstable case}
\label{sec:rel96-unstable}

For two uniform fluids the master equation in each region ($z\gtrless 0$)
admits solutions $e^{\pm kz}$ and $e^{\pm qz}$ with the relativistic
replacement
\begin{equation}
  q_{j}^{2}
    = \frac{n^{2}}{c^{2}}\;\frac{\enthalpy_{j}\,(c^{2}-\vA^{2})}{\vA^{2}\,\enthalpy_{j}}
    = \frac{n^{2}}{\vA^{2}}\;\Bigl(1 - \frac{\vA^{2}}{c^{2}}\Bigr)
    \quad(j=1,2),
  \label{eq:rel96-qj-gen}
\end{equation}
where $\vA$ is defined with the combined enthalpy.  More precisely,
defining
\begin{equation}
  \vA^{2}
    = \frac{b^{2}\,c^{2}}{\tfrac{1}{2}(\enthalpy_{1}+\enthalpy_{2}) + b^{2}},
  \label{eq:rel96-vA-combined}
\end{equation}
and the enthalpy fractions
\begin{equation}
  \alpha_{1,\text{rel}}
    = \frac{\enthalpy_{1}}{\enthalpy_{1}+\enthalpy_{2}},
  \qquad
  \alpha_{2,\text{rel}}
    = \frac{\enthalpy_{2}}{\enthalpy_{1}+\enthalpy_{2}},
  \label{eq:rel96-alpha-rel}
\end{equation}
the decay exponents become
\begin{equation}
  q_{j}
    = \frac{n}{\vA}\,\sqrt{\alpha_{j,\text{rel}}}\;
      \sqrt{1 - \frac{\vA^{2}}{c^{2}}}
  \qquad(n>0,\;j=1,2).
  \label{eq:rel96-qj}
\end{equation}
The factor $\sqrt{1-\vA^{2}/c^{2}}$ is the hallmark relativistic
correction: it shrinks the spatial decay length, concentrating the
perturbation closer to the interface as $\vA\to c$.

\paragraph{Boundary conditions.}
The continuity of $w$, $Dw$, and $D^{2}w$ across $z=0$ (cf.\
equations~\eqref{eq:10-187}) remains valid in the relativistic theory,
as does the jump condition obtained by integrating the master
equation~\eqref{eq:rel96-master} across the interface.  The algebra
proceeds exactly as in Chandrasekhar's derivation of the secular
equation~\eqref{eq:10-200}, with every occurrence of $\rho_{j}$
replaced by $\enthalpy_{j}/c^{2}$ and every $V_{a}$ replaced by
the relativistic $\vA$.

\paragraph{Relativistic dispersion relation.}
Non-dimensionalising with the units
\begin{equation}
  [n] = g/\vA,\qquad [k] = g/\vA^{2},
  \label{eq:rel96-units}
\end{equation}
the dispersion relation takes the form
\begin{equation}
  \boxed{%
  \hat{n}^{3}
    + 2\hat{k}\bigl(\sqrt{\alpha_{1,\text{rel}}}
      +\sqrt{\alpha_{2,\text{rel}}}\bigr)\,
      \Bigl(1-\frac{\vA^{2}}{c^{2}}\Bigr)^{1/2}\,\hat{n}^{2}
    + \hat{k}\Bigl[2\hat{k}\Bigl(1-\frac{\vA^{2}}{c^{2}}\Bigr)
      +\alpha_{1,\text{rel}}-\alpha_{2,\text{rel}}\Bigr]\hat{n}
    + 2\hat{k}^{3}\bigl(\sqrt{\alpha_{2,\text{rel}}}
      -\sqrt{\alpha_{1,\text{rel}}}\bigr)\,
      \Bigl(1-\frac{\vA^{2}}{c^{2}}\Bigr)^{1/2}
    = 0,}
  \label{eq:rel96-dispersion}
\end{equation}
where $\hat{n}=n/[n]$ and $\hat{k}=k/[k]$.  When $\vA/c\to0$, the
factor $(1-\vA^{2}/c^{2})^{1/2}\to1$ and
equation~\eqref{eq:rel96-dispersion} reduces identically to
Chandrasekhar's equation~\eqref{eq:10-209}.

\paragraph{Asymptotic limits.}
From~\eqref{eq:rel96-dispersion}:
\begin{equation}
  \left.
  \begin{aligned}
    n^{2}&\to gk\,(\alpha_{2,\text{rel}}-\alpha_{1,\text{rel}})
      &&(k\to0),\\
    n &\to \frac{g}{\vA}\,
      \bigl(\sqrt{\alpha_{2,\text{rel}}}-\sqrt{\alpha_{1,\text{rel}}}\bigr)\,
      \Bigl(1-\frac{\vA^{2}}{c^{2}}\Bigr)^{1/2}
      &&(k\to\infty).
  \end{aligned}
  \right\}
  \label{eq:rel96-asymptotics}
\end{equation}
The long-wavelength limit ($k\to0$) is \emph{unaffected} by the
magnetic field, exactly as in the classical theory.  However, the
short-wavelength saturation of the growth rate is reduced by the
factor $(1-\vA^{2}/c^{2})^{1/2}$: as $\vA\to c$, the maximum growth
rate is strongly suppressed.  This suppression reflects the additional
inertia contributed by the magnetic-field energy density in the
relativistic regime.

\paragraph{Physical interpretation.}
The factor $(1-\vA^{2}/c^{2})$ arises because, in the relativistic
theory, the magnetic tension responsible for stabilisation acts on the
total inertia $(\enthalpy+b^{2})/c^{2}$ rather than the mass
density~$\rho$ alone.  As the field strength increases, the Alfv\'en
speed approaches~$c$ but can never reach it; consequently, the magnetic
stabilisation, while enhanced, always remains \emph{causal}.

\subsection{Two uniform fluids: the stable case}
\label{sec:rel96-stable}

When $\alpha_{1,\text{rel}}>\alpha_{2,\text{rel}}$ (lighter fluid above
denser), the constant term in the cubic~\eqref{eq:rel96-dispersion}
is positive.  The Hurwitz criterion applied to the relativistic
cubic yields $BC-D>0$, $D>0$ (the algebra is identical to
equations~\eqref{eq:10-211}--\eqref{eq:10-212} with the replacement
$\alpha_{j}\to\alpha_{j,\text{rel}}$ and multiplication by
$(1-\vA^{2}/c^{2})^{1/2}$ in the appropriate places).  Hence all roots
have negative real parts, and the configuration is stable.

Setting $n=i\sigma$ for the oscillatory solutions, the perturbation in
each half-space takes the form
\begin{align}
  w_{1} &= A_{1}\,e^{+kz}
    + B_{1}^{-}\,e^{+i\sigma m_{1,\text{rel}}\,z}
    + C_{1}^{-}\,e^{-i\sigma m_{1,\text{rel}}\,z}
  \quad(z<0),
  \label{eq:rel96-stable-w1}\\
  w_{2} &= A_{2}\,e^{-kz}
    + B_{2}^{+}\,e^{-i\sigma m_{2,\text{rel}}\,z}
    + C_{2}^{+}\,e^{+i\sigma m_{2,\text{rel}}\,z}
  \quad(z>0),
  \label{eq:rel96-stable-w2}
\end{align}
where
\begin{equation}
  m_{j,\text{rel}}
    = \frac{1}{\vA}\,\sqrt{\alpha_{j,\text{rel}}}\;
      \sqrt{1-\frac{\vA^{2}}{c^{2}}}
  \label{eq:rel96-mj}
\end{equation}
is the relativistic wavenumber of the Alfv\'en waves in medium~$j$.
The oscillatory terms represent incoming and outgoing Alfv\'en waves
exactly as in the classical theory (equations~\eqref{eq:10-213}
and~\eqref{eq:10-214}), but now propagating at the relativistic Alfv\'en
speed~$\vA<c$.

As in the classical analysis, there is \emph{no dispersion relation
between~$k$ and~$\sigma$}: the boundary conditions at $z=0$ yield four
equations for six unknowns, leaving two free parameters that specify the
amplitudes and phases of the incoming Alfv\'en waves from $z=\pm\infty$.
The relativistic correction modifies only the propagation speed (and
hence the wavelength) of these Alfv\'en waves, not the qualitative
structure of the stable solution.

\paragraph{Causality.}
The Alfv\'en waves carry information at speed~$\vA<c$.  Since
$m_{j,\text{rel}}$ is real and finite for $\vA<c$, the wave solutions
are well-defined and causal.  In the ultra-relativistic limit
$\vA\to c$, the wavelength diverges ($m_{j,\text{rel}}\to0$) and the
oscillatory Alfv\'en modes decouple from the interface, leaving only the
evanescent solutions $e^{\pm kz}$.

%======================================================================
\section{The effect of a horizontal magnetic field --- relativistic
  extension}\label{sec:rel-97}
%======================================================================

We now extend \S97, where the impressed uniform field
$\mathbf{B}=B\hat{\mathbf{x}}$ is perpendicular to the gravitational
acceleration $\mathbf{g}=-g\hat{\mathbf{z}}$.

\subsection{Relativistic perturbation equations with horizontal field}

With the field along the $x$-axis, the linearised relativistic MHD
equations for perturbations
$\propto\exp\!\bigl[i(k_{x}x+k_{y}y)+nt\bigr]$ become:
\begin{gather}
  \frac{\enthalpy}{c^{2}}\,n\,u
    = -ik_{x}\,\delta p,
  \label{eq:rel97-x}\\
  \frac{\enthalpy}{c^{2}}\,n\,v
    - \frac{b^{2}}{c^{2}}\,\frac{ik_{x}}{n}
      \bigl(ik_{x}v-ik_{y}u\bigr)
    = -ik_{y}\,\delta p,
  \label{eq:rel97-y}\\
  \frac{\enthalpy}{c^{2}}\,n\,w
    - \frac{b^{2}}{c^{2}}\,\frac{ik_{x}}{n}
      \bigl(ik_{x}w-Du\bigr)
    = -D\,\delta p
      + \frac{g}{n}\,D\!\Bigl(\frac{\enthalpy}{c^{2}}\Bigr)\,w,
  \label{eq:rel97-z}\\
  \mathbf{h} = \frac{ik_{x}}{n}\,B_{0}\,\mathbf{u}.
  \label{eq:rel97-h}
\end{gather}
These are the relativistic analogues of
equations~\eqref{eq:10-215}--\eqref{eq:10-218}.

\subsection{Vanishing of the vorticity component $\zeta$}

Following the same combination of~\eqref{eq:rel97-x}
and~\eqref{eq:rel97-y} as in the classical theory, one obtains
\begin{equation}
  \Bigl(\frac{\enthalpy}{c^{2}}\,n
    + \frac{k_{x}^{2}}{n}\,\frac{b^{2}}{c^{2}}\Bigr)\,\zeta = 0,
  \label{eq:rel97-zeta-eq}
\end{equation}
where $\zeta = ik_{x}v - ik_{y}u$.  Since the coefficient is strictly
positive (both $\enthalpy>0$ and $b^{2}>0$), we conclude
\begin{equation}
  \zeta = 0.
  \label{eq:rel97-zeta-zero}
\end{equation}
This is the relativistic analogue of equation~\eqref{eq:10-223}.

\subsection{Relativistic master equation}

Eliminating $\delta p$ between the simplified horizontal and vertical
equations (following the algebraic route of
\S97), one obtains the relativistic master equation for a horizontal
field:
\begin{equation}
  D\!\Bigl(\frac{\enthalpy}{c^{2}}\,Dw\Bigr)
    + \frac{\vA^{2}}{n^{2}}\,\frac{\enthalpy}{c^{2}-\vA^{2}}
      \,\frac{k_{x}^{2}}{k^{2}}
      \,(D^{2}-k^{2})\,Dw
    - k^{2}\frac{\enthalpy}{c^{2}}\,w
    = -\frac{gk^{2}}{n^{2}}\,
      D\!\Bigl(\frac{\enthalpy}{c^{2}}\Bigr)\,w.
  \label{eq:rel97-master}
\end{equation}
This reduces to Chandrasekhar's equation~\eqref{eq:10-229} in the limit
$c\to\infty$.

\paragraph{Anisotropy.}
The magnetic term contains the factor $k_{x}^{2}/k^{2}=\cos^{2}\vartheta$,
where $\vartheta$ is the angle between the wave vector and the field
direction.  Consequently, the stabilising influence of the horizontal
field is \emph{anisotropic}: it is maximal for perturbations propagating
along the field ($\vartheta=0$) and vanishes for perturbations
perpendicular to it ($\vartheta=\pi/2$).  This anisotropy is inherited
directly from the classical theory and is not altered qualitatively by
relativistic corrections.

\subsection{Two uniform fluids: relativistic dispersion relation}

For two uniform fluids with $w\propto e^{\mp kz}$ in $z\gtrless 0$
(the solutions $e^{\pm kz}$ continue to hold since the magnetic
contribution in~\eqref{eq:rel97-master} does not alter the spatial
equation in homogeneous regions when $D(\enthalpy/c^{2})=0$), the
interface jump condition gives
\begin{equation}
  \boxed{%
  n^{2}
    = gk\,\biggl[
        \frac{\enthalpy_{2}-\enthalpy_{1}}{\enthalpy_{2}+\enthalpy_{1}}
        - \frac{2b^{2}\,k_{x}^{2}}{(\enthalpy_{2}+\enthalpy_{1})\,k}
          \;\frac{1}{1-\vA^{2}/c^{2}}
      \biggr].}
  \label{eq:rel97-dispersion}
\end{equation}
In the non-relativistic limit ($\enthalpy_{j}\to\rho_{j}c^{2}$,
$b^{2}=\mu H^{2}/(4\pi)$, $\vA^{2}/c^{2}\to0$), this reduces exactly
to Chandrasekhar's equation~\eqref{eq:10-234}:
\begin{equation}
  n^{2}\;\longrightarrow\;
  gk\,\biggl[\frac{\rho_{2}-\rho_{1}}{\rho_{2}+\rho_{1}}
    - \frac{\mu H^{2}k_{x}^{2}}{2\pi(\rho_{2}+\rho_{1})\,gk}\biggr]
  \quad(c\to\infty).
  \label{eq:rel97-NR-limit}
\end{equation}

\paragraph{Relativistic equivalent surface tension.}
By analogy with equation~\eqref{eq:10-235}, the magnetic field acts as
an equivalent surface tension
\begin{equation}
  T_{\text{eq,rel}}
    = \frac{b^{2}}{k}\,\frac{\cos^{2}\vartheta}{1-\vA^{2}/c^{2}}.
  \label{eq:rel97-Teq}
\end{equation}
The relativistic correction factor $1/(1-\vA^{2}/c^{2})$ \emph{enhances}
the equivalent surface tension compared with the classical value.  This
is physically intuitive: in the relativistic regime, the magnetic
tension $b^{2}$ acts on the enhanced inertia
$(\enthalpy+b^{2})/c^{2}$, and the net stabilising effect, expressed per
unit Atwood number, is amplified.

\subsection{Stability criterion and interchange instability}

From~\eqref{eq:rel97-dispersion}, instability ($n^{2}>0$) occurs when
the destabilising density contrast overcomes the magnetic stabilisation.
Setting $n^{2}=0$ determines the critical wavenumber:
\begin{equation}
  k_{x,\text{crit}}^{2}
    = \frac{(\enthalpy_{2}-\enthalpy_{1})\,gk}{2b^{2}}
      \;\bigl(1-\vA^{2}/c^{2}\bigr).
  \label{eq:rel97-kcrit}
\end{equation}
Modes with $k_{x}<k_{x,\text{crit}}$ are unstable; those with
$k_{x}>k_{x,\text{crit}}$ are stabilised by magnetic tension.

\paragraph{Interchange instability.}
When $k_{x}=0$ (perturbations perpendicular to $\mathbf{B}$), the
magnetic term in~\eqref{eq:rel97-dispersion} vanishes identically and
the dispersion relation reduces to the purely hydrodynamic form:
\begin{equation}
  n^{2}\big|_{k_{x}=0}
    = gk\,\frac{\enthalpy_{2}-\enthalpy_{1}}{\enthalpy_{2}+\enthalpy_{1}}.
  \label{eq:rel97-interchange}
\end{equation}
This is the \emph{relativistic interchange instability}: perturbations
that bend no field lines are unaffected by the magnetic field.  The only
relativistic modification is the replacement of the density ratio by the
enthalpy ratio, i.e.\ the Atwood number becomes
\begin{equation}
  \mathcal{A}_{\text{rel}}
    = \frac{\enthalpy_{2}-\enthalpy_{1}}{\enthalpy_{2}+\enthalpy_{1}}
  \label{eq:rel97-Atwood}
\end{equation}
rather than $(\rho_{2}-\rho_{1})/(\rho_{2}+\rho_{1})$.  For hot,
relativistic fluids (where $p\sim\varepsilon/3$), the enthalpy
$\enthalpy=\varepsilon+p=(4/3)\varepsilon$ differs substantially from
$\rho c^{2}$, and the relativistic Atwood number can differ markedly
from its classical counterpart.

%======================================================================
\section{Causality summary}\label{sec:rel-96-97-causality}
%======================================================================

Throughout the analysis of \S\S\ref{sec:rel-96} and~\ref{sec:rel-97},
causality is maintained by the fundamental bound
\begin{equation}
  \vA^{2}
    = \frac{b^{2}\,c^{2}}{\enthalpy + b^{2}}
    < c^{2}.
  \label{eq:rel96-97-causality}
\end{equation}
This bound ensures:
\begin{enumerate}
  \item The spatial decay exponents $q_{j}$ in the vertical-field
    problem (\S\ref{sec:rel-96}) are real and positive for unstable modes.
  \item The Alfv\'en wavenumbers $m_{j,\text{rel}}$ for stable modes
    are real, corresponding to finite-speed wave propagation.
  \item The equivalent surface tension $T_{\text{eq,rel}}$ in the
    horizontal-field problem (\S\ref{sec:rel-97}) remains finite.
  \item In the ultra-relativistic limit $\vA\to c$, growth rates are
    \emph{suppressed} (not enhanced), reflecting the increasing
    magnetic-field inertia.
\end{enumerate}
No additional constraint beyond $\vA<c$ is needed to ensure a causal
and well-posed relativistic magnetised RT instability theory.
